
\subsubsection{6.1 Interpretation}

The DNSI-gap correlation (R = -0.950) suggests that improvement entropy, normalized by task difficulty, may capture benchmark evolution dynamics. Higher DNSI (saturated benchmarks like CIFAR-10) corresponds to lower generalization gaps, while lower DNSI (less saturated benchmarks like HANS) corresponds to higher gaps.

The temporal prediction result ($R^2$ = 0.349) suggests potential as a leading indicator. Pre-2019 DNSI values show relationship with post-2019 gap measurements, with negative slope confirming the expected direction.

\subsubsection{6.2 Limitations}

\textbf{Small Sample Size (n = 4):} Bootstrap confidence intervals span the full range [-1, 1]. The p-value (0.050) is at the boundary of conventional significance. These results should be framed as a pilot study.

\textbf{Synthetic SOTA Histories:} The proof-of-concept uses historically-accurate synthetic data. Production implementation should use real PapersWithCode data.

\textbf{NLP Difficulty Proxy:} Class count maps naturally to vision tasks but less so to NLP tasks. The NLP correlation (R = -0.684) is based on estimated DNSI values. Domain-specific normalizers are needed for robust NLP application.

\textbf{Correlation $\neq$ Causation:} The results demonstrate predictive correlation, not causal mechanism. Saturation and overfitting may be confounded by third factors such as benchmark design quality.

\textbf{LOO-CV Instability:} Leave-one-out cross-validation yields $R^2$ = -3.82 for temporal prediction, indicating the model does not yet generalize reliably with n = 4.

\subsubsection{6.3 Scope Conditions}

Results apply to ML benchmarks with >15 SOTA entries over >3 years, where a difficulty proxy is available (class count for vision), and where held-out test sets exist for gap measurement. Results may not apply to NLP tasks without equivalent difficulty proxies, sparse or new benchmarks, or benchmarks without replication studies.

\subsubsection{6.4 Broader Implications}

DNSI could help researchers assess benchmark health using only historical SOTA records, without constructing expensive held-out test sets. The metric could potentially be gamed by artificially diversifying improvement patterns; DNSI should be one input among many for benchmark assessment.
